%================================================================
% CHAPTER 3: SOFTWARE REQUIREMENT SPECIFICATION
%================================================================
\chapter{SOFTWARE REQUIREMENT SPECIFICATION}

\section{Requirements Specifications}

In this chapter we describe the functional requirements and the non-functional requirements of the system briefly. Functional requirements are such type of requirements that are essential for the system; without these requirements the system is not considered a successful system. Functional requirements are those requirements which are essential for the system to perform a desired task or achieve a desired goal. In this case the user selects a university query, e.g., fee structure or other information of that particular university, then the crawler extracts that particular data from the university website and displays the result of the action performed by the user. While non-functional requirements are such type of requirements that specify criteria used to determine the operation of a device in preference to unique behaviors; they are contrasted with functional requirements that define specific behavior.

\section{Functional Requirements}

Functional requirements are such type of requirements that are essential for the system; without these requirements the system is not considered a successful system. Function requirements for the proposed system are listed below in Table~\ref{tab:basic_req}.

\subsection{Basic Requirements}

\begin{table}[H]
\centering
\caption{Basic Requirements}
\label{tab:basic_req}
\begin{tabularx}{\textwidth}{|c|X|}
\hline
\textbf{Sr. No.} & \textbf{Description} \\
\hline
R1.0 & This project is web based. \\
\hline
R1.1 & System would be responsive, used in computers and Android mobiles. \\
\hline
R1.2 & This project contains one menu and many different fields. \\
\hline
R1.3 & This project contains different master and web pages. \\
\hline
\end{tabularx}
\end{table}

\subsection{Searching}

In this module the user should be able to search the different universities of twin cities and also should be able to search for the entire program offered in which university. The user can select the city and program in which they want to take admission, then the system should display the list of those universities which offer this program as described in Table~\ref{tab:search_req}.

\begin{table}[H]
\centering
\caption{Search}
\label{tab:search_req}
\begin{tabularx}{\textwidth}{|c|X|}
\hline
\textbf{Sr. No.} & \textbf{Description} \\
\hline
R1.0 & The system should search universities of twin cities. \\
\hline
R1.1 & The system should search entire programs offered in which university. \\
\hline
\end{tabularx}
\end{table}

\subsection{Search General Information}

In this module the user should be able to view all the general information of different universities including fee structure, department, program offered, faculty, eligibility criteria, and admission date on one platform as described in Table~\ref{tab:general_req}.

\begin{table}[H]
\centering
\caption{General Information}
\label{tab:general_req}
\begin{tabularx}{\textwidth}{|c|X|}
\hline
\textbf{Sr. No.} & \textbf{Description} \\
\hline
R1.0 & User can view fee structure of all the programs. \\
\hline
R1.1 & It contains also complete information about faculties of different departments. \\
\hline
R1.2 & User can easily get all information about departments. \\
\hline
R1.3 & User can easily view which programs are offered in which university. \\
\hline
R1.4 & User can easily get the information about all the universities' admission open and close dates. \\
\hline
R1.5 & User can easily see the eligibility criteria of different courses of different universities. \\
\hline
\end{tabularx}
\end{table}

\subsection{Comparing Information}

In this module the user should be able to compare universities on different criteria, e.g., fee structure and university ranking. Users compare fee structure by selecting program and universities. The system shows the fee of the entire program of all the universities, then users can easily see the comparison without visiting all the universities one by one. Users also compare universities on the basis of HEC ranking of the university or entire program, as given in Table~\ref{tab:comparison_req}.

\begin{table}[H]
\centering
\caption{Comparison}
\label{tab:comparison_req}
\begin{tabularx}{\textwidth}{|c|X|}
\hline
\textbf{Sr. No.} & \textbf{Description} \\
\hline
R1.0 & User can easily view the comparison of fee structure of different university programs. \\
\hline
R1.1 & User can easily view the comparison of different university programs on HEC ranking. \\
\hline
\end{tabularx}
\end{table}

\section{Non-Functional Requirements}

Non-functional requirements specify criteria that are used to determine the operation of a device in preference to unique behaviors. All these features help the functional requirements to achieve the goal. Non-functional requirements include efficiency, effectiveness, reliability, resources, project design and availability.

\subsection{Interface Design}

The Academia Portal design is easy to use and navigate for all information. This design is created for users so they can easily and efficiently search the universities, and the user of this website can search all the information in minimum time and effort. The design of this project is simple for users through which they can easily navigate and explore icons.

\subsection{Reliability}

The particularly non-functional requirement of the existing project is reliability. We are quite assured that the proposed system gives a great level of accuracy. The system assures the availability of the right information at the right time. All the information is precise, safe and operated properly.

\subsection{Efficiency}

The proposed system Academia Portal is more effective in all ways as compared to the existing system. The proposed system extracts data from different websites and integrates them on one site, which is not available in any existing system.

\subsection{Availability}

Availability is also a very important factor. Availability of the system is 24 hours seven days a week, which means that the system must be available at any time for users. Another important factor is that the system should be available online.

\subsection{Performance}

The current system performs each function efficiently in less time. After testing the proposed system, the processing time, delivery time and its response to users is satisfactory and efficient.

\subsection{Accuracy}

Accuracy generally defines the level of similarity between the values which are known and which are taken. It gives the way to show how much the project is workable. Functionality should be performed in an accurate way; in case of failure of any module, other modules will not be affected.

\subsection{Ease of Use}

This project is web based. The advanced Academia web portal provides a user-friendly environment for users which is easy to use and understandable. It is best for all categories of users: beginner, novice and expert.

\subsection{User Friendly}

The system should be user friendly and give an easy way to understand the system. This system makes the environment friendly and is very easy for all, and smart in the order to accomplish the responsibility of understanding.
